% Auto-generated from raw.txt
% Usage:
%   \vocabentry{word}{/ipa/ (pos)}{definition}{example}{note}

\vocabentry{Macroeconomics}
{/ˌmækroʊˌekəˈnɑːmɪks/ (n)}
{The part of economics concerned with large-scale or general economic factors.}
{Macroeconomics helps governments understand how to control inflation and unemployment rates.}
{Đây là một danh từ dùng để chỉ kinh tế học vĩ mô, nghiên cứu các yếu tố tổng thể như GDP hay lạm phát.}

\vocabentry{Microeconomics}
{/ˌmaɪkroʊˌekəˈnɑːmɪks/ (n)}
{The part of economics concerned with single factors and the effects of individual decisions.}
{A course in microeconomics explores how supply and demand determine prices in specific markets.}
{Đây là một danh từ dùng để chỉ kinh tế học vi mô, nghiên cứu hành vi của cá nhân và doanh nghiệp.}

\vocabentry{Inflation}
{/ɪnˈfleɪʃn/ (n)}
{A general increase in prices and fall in the purchasing value of money.}
{High inflation can erode the savings of many citizens over time.}
{Đây là một danh từ dùng để chỉ sự lạm phát, hiện tượng mức giá chung tăng lên.}

\vocabentry{Deflation}
{/ˌdiːˈfleɪʃn/ (n)}
{Reduction of the general level of prices in an economy.}
{Persistent deflation can lead to a decrease in consumer spending as people wait for lower prices.}
{Đây là một danh từ dùng để chỉ sự giảm phát, mức giá chung của hàng hóa và dịch vụ giảm xuống.}

\vocabentry{Fiscal policy}
{/ˈfɪskl ˈpɑːləsi/ (n)}
{Government policy that attempts to influence the direction of the economy through changes in government spending or taxes.}
{The government implemented a new fiscal policy to stimulate economic growth during the recession.}
{Đây là một cụm danh từ dùng để chỉ chính sách tài khóa.}

\vocabentry{Monetary policy}
{/ˈmɑːnɪteri ˈpɑːləsi/ (n)}
{Policy that involves altering the quantity of money in circulation and the interest rate.}
{The Central Bank adjusted its monetary policy to curb rising inflation.}
{Đây là một cụm danh từ dùng để chỉ chính sách tiền tệ do ngân hàng trung ương điều hành.}

\vocabentry{Gross Domestic Product}
{/ɡroʊs dəˈmestɪk ˈprɑːdʌkt/ (n)}
{The total value of goods produced and services provided in a country during one year.}
{GDP is a primary indicator used to gauge the health of a country's economy.}
{Đây là một cụm danh từ dùng để chỉ tổng sản phẩm quốc nội - GDP.}

\vocabentry{Recession}
{/rɪˈseʃn/ (n)}
{A period of temporary economic decline during which trade and industrial activity are reduced.}
{The country fell into a deep recession following the global financial crisis.}
{Đây là một danh từ dùng để chỉ sự suy thoái kinh tế.}

\vocabentry{Subsidy}
{/ˈsʌbsədi/ (n)}
{A sum of money granted by the government or a public body to assist an industry or business.}
{Agricultural subsidies help farmers stay competitive in the international market.}
{Đây là một danh từ dùng để chỉ tiền trợ cấp hoặc bao cấp của chính phủ.}

\vocabentry{Revenue}
{/ˈrevənuː/ (n)}
{Income, especially when of a company or organization and of a substantial nature.}
{The company reported a significant increase in revenue thanks to its new product line.}
{Đây là một danh từ dùng để chỉ doanh thu hoặc nguồn thu ngân sách.}

\vocabentry{Expenditure}
{/ɪkˈspendɪtʃər/ (n)}
{An amount of money spent, as by a government or person.}
{The government's expenditure on healthcare has doubled over the last decade.}
{Đây là một danh từ dùng để chỉ sự chi tiêu hoặc khoản chi.}

\vocabentry{Deficit}
{/ˈdefɪsɪt/ (n)}
{The amount by which something, especially a sum of money, is too small.}
{The budget deficit has forced the government to cut spending on social programs.}
{Đây là một danh từ dùng để chỉ sự thâm hụt, thường dùng cho thâm hụt ngân sách.}

\vocabentry{Surplus}
{/ˈsɜːrpləs/ (n)}
{An amount of something left over when requirements have been met.}
{The trade surplus indicates that the country is exporting more than it is importing.}
{Đây là một danh từ dùng để chỉ sự dư thừa hoặc thặng dư ngân sách.}

\vocabentry{Liabilities}
{/ˌlaɪəˈbɪlətiz/ (n)}
{A thing for which someone is responsible, especially a debt or financial obligation.}
{The company's liabilities exceed its assets, putting it at risk of bankruptcy.}
{Đây là một danh từ dùng để chỉ các khoản nợ phải trả hoặc trách nhiệm pháp lý.}

\vocabentry{Assets}
{/ˈæsets/ (n)}
{A useful or valuable thing, person, or quality.}
{Real estate and stocks are considered fixed and liquid assets, respectively.}
{Trong kinh tế, đây là danh từ dùng để chỉ tài sản thuộc sở hữu của một cá nhân hay doanh nghiệp.}

\vocabentry{Liquidity}
{/lɪˈkwɪdəti/ (n)}
{The availability of liquid assets to a market or company.}
{High liquidity is important for businesses to meet their short-term financial obligations.}
{Đây là một danh từ dùng để chỉ tính thanh khoản, khả năng chuyển đổi tài sản thành tiền mặt.}

\vocabentry{Equity}
{/ˈekwəti/ (n)}
{The value of the shares issued by a company.}
{He decided to invest in the company's equity to participate in its future profits.}
{Đây là một danh từ dùng để chỉ vốn chủ sở hữu hoặc giá trị cổ phần.}

\vocabentry{Dividend}
{/ˈdɪvɪdend/ (n)}
{A sum of money paid regularly by a company to its shareholders out of its profits.}
{Investors were pleased when the company announced an increase in annual dividends.}
{Đây là một danh từ dùng để chỉ cổ tức.}

\vocabentry{Capital}
{/ˈkæpɪtl/ (n)}
{Wealth in the form of money or other assets owned by a person or organization.}
{Small businesses often struggle to raise the initial capital needed to start operations.}
{Đây là một danh từ dùng để chỉ vốn đầu tư.}

\vocabentry{Market capitalization}
{/ˈmɑːrkɪt ˌkæpɪtələˈzeɪʃn/ (n)}
{The total value of a company's shares on the stock market.}
{Apple was the first company to reach a market capitalization of three trillion dollars.}
{Đây là một cụm danh từ dùng để chỉ giá trị vốn hóa thị trường.}

\vocabentry{Commodity}
{/kəˈmɑːdəti/ (n)}
{A raw material or primary agricultural product that can be bought and sold.}
{Gold and oil are two of the most heavily traded commodities in the world.}
{Đây là một danh từ dùng để chỉ hàng hóa cơ bản hoặc nguyên liệu thô.}

\vocabentry{Privatization}
{/ˌpraɪvətəˈzeɪʃn/ (n)}
{The transfer of a business, industry, or service from public to private ownership and control.}
{The privatization of the water supply remains a highly controversial issue.}
{Đây là một danh từ dùng để chỉ sự tư nhân hóa các dịch vụ công.}

\vocabentry{Monopoly}
{/əˈmɑːnəpəli/ (n)}
{The exclusive possession or control of the supply of or trade in a commodity or service.}
{The government is taking steps to prevent a single company from holding a monopoly on internet services.}
{Đây là một danh từ dùng để chỉ sự độc quyền.}

\vocabentry{Oligopoly}
{/ˌɑːlɪˈɡɑːpəli/ (n)}
{A state of limited competition, in which a market is shared by a small number of producers or sellers.}
{The smartphone industry is largely an oligopoly dominated by a few major players.}
{Đây là một danh từ dùng để chỉ thị trường độc quyền tập đoàn.}

\vocabentry{Diversification}
{/daɪˌvɜːrsɪfɪˈkeɪʃn/ (n)}
{The process of a business enlarging or varying its range of products or field of operation.}
{Portfolio diversification is a key strategy for managing financial risk.}
{Đây là một danh từ dùng để chỉ sự đa dạng hóa, đặc biệt là đa dạng hóa danh mục đầu tư.}

\vocabentry{Interest rate}
{/ˈɪntrest reɪt/ (n)}
{The proportion of a loan that is charged as interest to the borrower.}
{Rising interest rates make borrowing money more expensive for consumers and businesses.}
{Đây là một cụm danh từ dùng để chỉ lãi suất.}

\vocabentry{Mortgage}
{/ˈmɔːrɡɪdʒ/ (n/v)}
{A legal agreement by which a bank lends money at interest in exchange for taking title of the debtor's property.}
{They took out a thirty-year mortgage to buy their first family home.}
{Đây là một danh từ/động từ dùng để chỉ sự thế chấp hoặc khoản vay thế chấp bất động sản.}

\vocabentry{Bankruptcy}
{/ˈbæŋkrʌptsi/ (n)}
{The state of being bankrupt.}
{The economic downturn led to a wave of corporate bankruptcies across the country.}
{Đây là một danh từ dùng để chỉ sự phá sản.}

\vocabentry{Credit rating}
{/ˈkredɪt ˈreɪtɪŋ/ (n)}
{An estimate of the ability of a person or organization to fulfill their financial commitments.}
{A poor credit rating can make it very difficult to get a loan or a credit card.}
{Đây là một cụm danh từ dùng để chỉ xếp hạng tín nhiệm.}

\vocabentry{Collateral}
{/kəˈlætərəl/ (n)}
{Something pledged as security for repayment of a loan, to be forfeited in the event of a default.}
{He used his house as collateral for the business loan.}
{Đây là một danh từ dùng để chỉ tài sản thế chấp.}

\vocabentry{Equilibrium}
{/ˌiːkwɪˈlɪbriəm/ (n)}
{A state in which opposing forces or influences are balanced.}
{Market equilibrium occurs when the quantity demanded equals the quantity supplied.}
{Trong kinh tế, đây là trạng thái cân bằng giữa cung và cầu.}

\vocabentry{Scarcity}
{/ˈskersəti/ (n)}
{The state of being scarce or in short supply; shortage.}
{Scarcity of raw materials has led to a sharp increase in production costs.}
{Đây là một danh từ dùng để chỉ sự khan hiếm tài nguyên.}

\vocabentry{Opportunity cost}
{/ˌɑːpərˈtuːnəti kɔːst/ (n)}
{The loss of potential gain from other alternatives when one alternative is chosen.}
{The opportunity cost of going to university is the wages you could have earned during those four years.}
{Đây là một cụm danh từ dùng để chỉ chi phí cơ hội.}

\vocabentry{Incentive}
{/ɪnˈsentɪv/ (n)}
{A thing that motivates or encourages someone to do something.}
{The government is offering tax incentives to companies that relocate to rural areas.}
{Đây là một danh từ dùng để chỉ sự khuyến khích hoặc động lực.}

\vocabentry{Utility}
{/juːˈtɪləti/ (n)}
{The state of being useful, profitable, or beneficial.}
{Marginal utility is the additional satisfaction a consumer gains from consuming one more unit of a good.}
{Trong kinh tế học, đây là danh từ dùng để chỉ độ thỏa dụng hoặc ích lợi của hàng hóa.}

\vocabentry{Elasticity}
{/ˌiːlæˈstɪsəti/ (n)}
{The degree to which a demand or supply is sensitive to changes in price or income.}
{The price elasticity of demand for luxury goods is usually very high.}
{Đây là một danh từ dùng để chỉ độ co giãn của cung hoặc cầu.}

\vocabentry{Externalities}
{/ˌekstɜːrˈnælətiz/ (n)}
{A side effect or consequence of an industrial or commercial activity that affects other parties.}
{Pollution is a negative externality that is often not included in the cost of production.}
{Đây là một danh từ dùng để chỉ ngoại ứng, có thể là tích cực hoặc tiêu cực.}

\vocabentry{Public goods}
{/ˌpʌblɪk ˈɡʊdz/ (n)}
{A commodity or service that is provided without profit to all members of a society.}
{Clean air and national defense are classic examples of public goods.}
{Đây là một cụm danh từ dùng để chỉ hàng hóa công cộng.}

\vocabentry{Information asymmetry}
{/ˌɪnfərˈmeɪʃn eɪˈsɪmətri/ (n)}
{A situation in which one party in a transaction has more or superior information compared to another.}
{Information asymmetry in the used car market can lead to buyers paying too much for poor-quality vehicles.}
{Đây là một cụm danh từ dùng để chỉ sự bất đối xứng thông tin.}

\vocabentry{Market failure}
{/ˈmɑːrkɪt ˈfeɪljər/ (n)}
{A situation in which the allocation of goods and services by a free market is not efficient.}
{Governments often intervene in the economy to correct market failures like monopolies.}
{Đây là một cụm danh từ dùng để chỉ sự thất bại của thị trường.}

\vocabentry{Exchange rate}
{/ɪksˈtʃeɪndʒ reɪt/ (n)}
{The value of one currency for the purpose of conversion to another.}
{A strong exchange rate can make a country's exports more expensive for foreign buyers.}
{Đây là một cụm danh từ dùng để chỉ tỷ giá hối đoái.}

\vocabentry{Balance of payments}
{/ˈbæləns əv ˈpeɪmənts/ (n)}
{The difference in total value between payments into and out of a country over a period.}
{A persistent deficit in the balance of payments can put pressure on a country's currency.}
{Đây là một cụm danh từ dùng để chỉ cán cân thanh toán.}

\vocabentry{Tariff}
{/ˈtærɪf/ (n)}
{A tax or duty to be paid on a particular class of imports or exports.}
{The government imposed new tariffs on imported steel to protect domestic manufacturers.}
{Đây là một danh từ dùng để chỉ thuế quan, một rào cản thương mại.}

\vocabentry{Quota}
{/ˈkwoʊtə/ (n)}
{A fixed share of something that a person or group is entitled to receive or is bound to contribute.}
{The import quota on textiles has been relaxed to encourage free trade.}
{Trong thương mại, đây là hạn ngạch nhập khẩu.}

\vocabentry{Embargo}
{/ɪmˈbɑːrɡoʊ/ (n)}
{An official ban on trade or other commercial activity with a particular country.}
{The international community imposed an oil embargo on the country in response to its human rights violations.}
{Đây là một danh từ dùng để chỉ lệnh cấm vận thương mại.}

\vocabentry{Sanction}
{/ˈsæŋkʃn/ (n/v)}
{A threatened penalty for disobeying a law or rule.}
{The lifting of international sanctions has allowed the country to start exporting oil again.}
{Đây là một danh từ/động từ dùng để chỉ lệnh trừng phạt hoặc việc xử phạt.}

\vocabentry{Sovereign debt}
{/ˈsɑːvrən det/ (n)}
{The amount of money that a country's government has borrowed, typically through issuing bonds.}
{High levels of sovereign debt can make it difficult for a country to attract investment.}
{Đây là một cụm danh từ dùng để chỉ nợ công hoặc nợ chủ quyền.}

\vocabentry{Cryptocurrency}
{/ˈkrɪptoʊkɜːrənsi/ (n)}
{A digital currency in which transactions are verified and records maintained by a decentralized system.}
{Bitcoin is the most well-known cryptocurrency in the global financial market.}
{Đây là một danh từ dùng để chỉ tiền mã hóa hoặc tiền điện tử.}

\vocabentry{Blockchain}
{/ˈblɑːktʃeɪn/ (n)}
{A system in which a record of transactions made in bitcoin or another cryptocurrency are maintained.}
{Blockchain technology has the potential to revolutionize how financial transactions are recorded and verified.}
{Đây là một danh từ dùng để chỉ công nghệ chuỗi khối.}

\vocabentry{Fintech}
{/ˈfɪntek/ (n)}
{Computer programs and other technology used to support or enable banking and financial services.}
{The rise of fintech has made it easier for people to manage their money and access loans.}
{Đây là một danh từ dùng để chỉ công nghệ tài chính.}

\vocabentry{Hedge fund}
{/ˈhedʒ fʌnd/ (n)}
{A limited partnership of investors that uses high-risk methods, such as investing with borrowed money, in hopes of realizing large capital gains.}
{He works as an analyst for a large hedge fund in London.}
{Đây là một danh từ dùng để chỉ quỹ đầu tư mạo hiểm.}

\vocabentry{Stock exchange}
{/ˈstɑːk ɪksˌtʃeɪndʒ/ (n)}
{A market in which securities are bought and sold.}
{The New York Stock Exchange is the largest in the world by market capitalization.}
{Đây là một danh từ dùng để chỉ sàn giao dịch chứng khoán.}

\vocabentry{Blue-chip}
{/ˌbluː ˈtʃɪp/ (adj)}
{Denoting companies or their shares considered to be a reliable investment.}
{Investors often prefer blue-chip stocks during periods of economic uncertainty.}
{Đây là một tính từ dùng để chỉ các cổ phiếu của công ty lớn, uy tín và ổn định.}

\vocabentry{Bull market}
{/ˈbʊl ˌmɑːrkɪt/ (n)}
{A market in which share prices are rising, encouraging buying.}
{The tech sector has been in a bull market for several years now.}
{Đây là một cụm danh từ dùng để chỉ thị trường bò tót - thị trường tăng trưởng.}

\vocabentry{Bear market}
{/ˈber ˌmɑːrkɪt/ (n)}
{A market in which prices are falling, encouraging selling.}
{The sudden bear market caused panic among many small investors.}
{Đây là một cụm danh từ dùng để chỉ thị trường gấu - thị trường suy thoái.}

\vocabentry{Portfolio}
{/pɔːrtˈfoʊlioʊ/ (n)}
{A range of investments held by a person or organization.}
{A well-balanced portfolio should include a mix of stocks, bonds, and cash.}
{Đây là một danh từ dùng để chỉ danh mục đầu tư.}

\vocabentry{Arbitrage}
{/ˈɑːrbɪtrɑːʒ/ (n)}
{The simultaneous buying and selling of securities, currency, or commodities in different markets.}
{Large financial institutions often use automated systems to exploit arbitrage opportunities.}
{Đây là một danh từ dùng để chỉ hoạt động kinh doanh chênh lệch giá.}

\vocabentry{Speculation}
{/ˌspekjuˈleɪʃn/ (n)}
{Investment in stocks, property, or other ventures in the hope of gain but with the risk of loss.}
{Excessive speculation in the housing market can lead to a dangerous bubble.}
{Đây là một danh từ dùng để chỉ hoạt động đầu cơ.}

\vocabentry{Insider trading}
{/ɪnˌsaɪdər ˈtreɪdɪŋ/ (n)}
{The illegal practice of trading on the stock exchange to one's own advantage through having access to confidential information.}
{The former CEO was arrested and charged with insider trading.}
{Đây là một cụm danh từ dùng để chỉ hành vi giao dịch nội gián.}

\vocabentry{Capital gain}
{/ˈkæpɪtl ɡeɪn/ (n)}
{A profit from the sale of property or an investment.}
{You have to pay tax on any capital gains you make from selling your shares.}
{Đây là một cụm danh từ dùng để chỉ lợi nhuận vốn hoặc lãi vốn.}

\vocabentry{Wealth tax}
{/welθ tæk/ (n)}
{A tax on the total value of assets owned by an individual.}
{Some economists argue that a wealth tax is necessary to reduce income inequality.}
{Đây là một cụm danh từ dùng để chỉ thuế tài sản.}

\vocabentry{Fiscal year}
{/ˈfɪskl jɪər/ (n)}
{A period used by a government or business for accounting purposes and preparing financial statements.}
{The company's fiscal year begins on April 1st and ends on March 31st.}
{Đây là một cụm danh từ dùng để chỉ năm tài chính.}

\vocabentry{Direct tax}
{/dəˌrekt ˈtæks/ (n)}
{A tax, such as income tax, which is levied on the income or profits of the person who pays it.}
{Income tax and corporate tax are two examples of direct taxes.}
{Đây là một cụm danh từ dùng để chỉ thuế trực thu.}

\vocabentry{Indirect tax}
{/ˌɪndəˈrekt ˈtæks/ (n)}
{A tax levied on goods or services rather than on persons or organizations.}
{Value Added Tax (VAT) is a common form of indirect tax in many countries.}
{Đây là một cụm danh từ dùng để chỉ thuế gián thu, ví dụ như VAT.}

\vocabentry{Tax evasion}
{/ˈtæks ɪˌveɪʒn/ (n)}
{The illegal non-payment or underpayment of tax.}
{The government is launching a crackdown on tax evasion by wealthy individuals.}
{Đây là một cụm danh từ dùng để chỉ hành vi trốn thuế.}

\vocabentry{Tax avoidance}
{/ˈtæks əˈvɔɪdns/ (n)}
{The arrangement of one's financial affairs to minimize tax liability within the law. (Đây là một cụm danh từ dùng để chỉ việc tránh thuế (hợp pháp nhưng gây tranh cãi).).}
{Many multinational corporations use legal tax avoidance strategies to reduce their tax bills.}
{}

\vocabentry{Brackets}
{/ˈbrækɪts/ (n)}
{Each of several groups of people with incomes within particular limits for the purpose of taxation.}
{He recently moved into a higher tax bracket after receiving a promotion.}
{Trong thuế, đây là danh từ dùng để chỉ các bậc thuế/khung thuế.}

\vocabentry{Auditing}
{/ˈɔːdɪtɪŋ/ (n)}
{The process of examining a person's or organization's accounts.}
{Regular auditing is essential for ensuring transparency and preventing fraud in large companies.}
{Đây là một danh từ dùng để chỉ hoạt động kiểm toán.}

\vocabentry{Consolidation}
{/kənˌsɑːlɪˈdeɪʃn/ (n)}
{The action or process of combining several things into a single more effective or coherent whole.}
{Debt consolidation allows you to combine multiple loans into a single monthly payment.}
{Trong kinh tế, đây là sự hợp nhất hoặc củng cố tài chính.}

\vocabentry{Liability}
{/ˌlaɪəˈbɪləti/ (n)}
{Legal responsibility for something; in finance, a company's debts or obligations.}
{The company has several major liabilities that it needs to settle before the end of the year.}
{Đây là một danh từ dùng để chỉ trách nhiệm pháp lý hoặc khoản nợ phải trả.}

\vocabentry{Amortization}
{/əˌmɔːrtəˈzeɪʃn/ (n)}
{The action or process of gradually writing off the initial cost of an asset.}
{The mortgage includes a detailed amortization schedule for the entire thirty-year period.}
{Đây là một danh từ dùng để chỉ sự khấu hao tài sản vô hình hoặc trả dần nợ.}

\vocabentry{Depreciation}
{/dɪˌpriːʃiˈeɪʃn/ (n)}
{A reduction in the value of an asset over time, due in particular to wear and tear.}
{The depreciation of the currency has made imports more expensive for local businesses.}
{Đây là một danh từ dùng để chỉ sự khấu hao tài sản hữu hình hoặc giảm giá trị tiền tệ.}

\vocabentry{Revaluation}
{/ˌriːˌvæljuˈeɪʃn/ (n)}
{An official change in the value of a country's currency relative to other currencies. (Đây là một danh từ dùng để chỉ sự định giá lại (thường là tăng giá) tiền tệ.).}
{The government announced a revaluation of the national currency to combat inflation.}
{}

\vocabentry{Appreciation}
{/əˌpriːʃiˈeɪʃn/ (n)}
{An increase in the value of an asset over time.}
{The appreciation of house prices in the city has made it difficult for first-time buyers.}
{Đây là một danh từ dùng để chỉ sự tăng giá trị tài sản.}

\vocabentry{Solvency}
{/ˈsɑːlvənsi/ (n)}
{The possession of assets in excess of liabilities; ability to pay one's debts.}
{The central bank is monitoring the solvency of several smaller banks during the crisis.}
{Đây là một danh từ dùng để chỉ khả năng thanh toán nợ lâu dài.}

\vocabentry{Insolvency}
{/ɪnˈsɑːlvənsi/ (n)}
{The state of being unable to pay one's debts.}
{The company was forced into insolvency after its main customer went bankrupt.}
{Đây là một danh từ dùng để chỉ tình trạng mất khả năng thanh toán nợ.}

\vocabentry{Venture capital}
{/ˈventʃər ˈkæpɪtl/ (n)}
{Capital invested in a project in which there is a substantial element of risk.}
{Silicon Valley is famous for its high concentration of venture capital firms.}
{Đây là một cụm danh từ dùng để chỉ vốn đầu tư mạo hiểm cho các startup.}

\vocabentry{Angel investor}
{/ˈeɪndʒl ɪnˌvestər/ (n)}
{A wealthy individual who provides capital for a business start-up, usually in exchange for convertible debt or ownership equity.}
{The startup was able to secure funding from an angel investor to develop its prototype.}
{Đây là một cụm danh từ dùng để chỉ nhà đầu tư thiên thần.}

\vocabentry{Crowdfunding}
{/ˈkraʊdfʌndɪŋ/ (n)}
{The practice of funding a project or venture by raising small amounts of money from a large number of people.}
{They used a crowdfunding platform to raise money for their new documentary film.}
{Đây là một danh từ dùng để chỉ việc gọi vốn cộng đồng.}

\vocabentry{Underwriting}
{/ˈʌndərraɪtɪŋ/ (n)}
{The process by which an individual or institution takes on financial risk for a fee.}
{Several large banks were involved in the underwriting of the company's initial public offering.}
{Đây là một danh từ dùng để chỉ hoạt động bảo lãnh phát hành chứng khoán hoặc bảo hiểm.}

\vocabentry{IPO}
{/ˌaɪ piː ˈoʊ/ (n)}
{Initial Public Offering; the first sale of stock by a private company to the public.}
{The company's IPO was a huge success, raising over five hundred million dollars.}
{Đây là một danh từ viết tắt dùng để chỉ sự phát hành cổ phiếu lần đầu ra công chúng.}

\vocabentry{Prospectus}
{/prəˈspektəs/ (n)}
{A printed document that advertises or describes a school, commercial enterprise, etc., in order to attract or inform clients or members.}
{You should read the investment prospectus carefully before buying any shares.}
{Trong kinh tế là bản cáo bạch.}

\vocabentry{Bullish}
{/ˈbʊlɪʃ/ (adj)}
{Characterized by rising share prices; or being optimistic about the future.}
{Analysts are bullish about the company's prospects for the next quarter.}
{Đây là một tính từ dùng để chỉ xu hướng lạc quan, tin rằng giá sẽ tăng.}

\vocabentry{Bearish}
{/ˈberɪʃ/ (adj)}
{Characterized by falling share prices; or being pessimistic about the future.}
{The sudden drop in oil prices has made investors bearish about energy stocks.}
{Đây là một tính từ dùng để chỉ xu hướng bi quan, tin rằng giá sẽ giảm.}

\vocabentry{Volatility}
{/ˌvɑːləˈtɪləti/ (n)}
{Liability to change rapidly and unpredictably, especially for the worse.}
{Market volatility can make it very difficult for investors to plan for the long term.}
{Đây là một danh từ dùng để chỉ sự biến động mạnh của thị trường.}

\vocabentry{Leverage}
{/ˈlevərɪdʒ/ (n/v)}
{The use of borrowed money to increase the potential return of an investment.}
{High levels of financial leverage can increase profits but also significantly increase risk.}
{Đây là một danh từ/động từ dùng để chỉ đòn bẩy tài chính.}

\vocabentry{Derivative}
{/dɪˈrɪvətɪv/ (n)}
{A financial instrument whose value is based on one or more underlying assets.}
{Options and futures are common types of derivatives used for hedging risk.}
{Đây là một danh từ dùng để chỉ công cụ tài chính phái sinh.}

\vocabentry{Hedging}
{/ˈhedʒɪŋ/ (n)}
{Protecting oneself against loss on an investment by making a balancing or compensating investment.}
{Gold is often used as a tool for hedging against inflation and economic uncertainty.}
{Đây là một danh từ dùng để chỉ hoạt động phòng vệ rủi ro.}

\vocabentry{Option}
{/ˈɑːpʃn/ (n)}
{A contract that gives the buyer the right, but not the obligation, to buy or sell an asset.}
{He bought a call option on the stock, hoping that the price would rise.}
{Đây là một danh từ dùng để chỉ quyền chọn trong đầu tư.}

\vocabentry{Future}
{/ˈfjuːtʃər/ (n)}
{A contract to buy or sell an asset at a predetermined price at a specified time in the future.}
{Farmers use futures contracts to lock in a price for their crops before they are harvested.}
{Đây là một danh từ dùng để chỉ hợp đồng tương lai.}

\vocabentry{Arbitrageur}
{/ˌɑːrbɪtrɑːˈʒɜːr/ (n)}
{A person who engages in arbitrage.}
{Arbitrageurs play an important role in ensuring that prices remain consistent across different markets.}
{Đây là một danh từ dùng để chỉ người thực hiện kinh doanh chênh lệch giá.}

\vocabentry{Mercantilism}
{/ˈmɜːrkəntɪlɪzəm/ (n)}
{An economic theory that trade generates wealth and is stimulated by the accumulation of profitable balances.}
{Mercantilism was the dominant economic theory in Europe during the 16th to 18th centuries.}
{Đây là một danh từ dùng để chỉ chủ nghĩa trọng thương.}

\vocabentry{Protectionism}
{/prəˈtekʃənɪzəm/ (n)}
{The theory or practice of shielding a country's domestic industries from foreign competition.}
{Critics argue that protectionism can lead to higher prices for consumers and slower economic growth.}
{Đây là một danh từ dùng để chỉ chủ nghĩa bảo hộ mậu dịch.}

\vocabentry{Laissez-faire}
{/ˌleseɪ ˈfer/ (n)}
{A policy or attitude of letting things take their own course, without interfering.}
{Laissez-faire economics advocates for minimal government intervention in the market.}
{Đây là một danh từ dùng để chỉ học thuyết kinh tế tự do không can thiệp.}

\vocabentry{Keynesianism}
{/ˌkeɪnziəˈnɪzəm/ (n)}
{The economic theories of John Maynard Keynes, especially those advocating government spending to stimulate demand.}
{Keynesianism became popular after the Great Depression as a way to manage economic cycles.}
{Đây là một danh từ dùng để chỉ thuyết kinh tế Keynes.}

\vocabentry{Monetarism}
{/ˈmʌnɪtərɪzəm/ (n)}
{The theory that inflation is caused by an excess of money in the economy.}
{Monetarism emphasizes the importance of controlling the money supply to maintain price stability.}
{Đây là một danh từ dùng để chỉ thuyết trọng tiền.}

\vocabentry{Stagflation}
{/ˌstæɡˈfleɪʃn/ (n)}
{Persistent high inflation combined with high unemployment and stagnant demand.}
{The country experienced a period of stagflation in the 1970s due to rising oil prices.}
{Đây là một danh từ dùng để chỉ tình trạng lạm phát đình trệ.}

\vocabentry{Hyperinflation}
{/ˌhaɪpərɪnˈfleɪʃn/ (n)}
{Monetary inflation occurring at a very high rate.}
{Hyperinflation in Germany during the 1920s led to a complete collapse of the national currency.}
{Đây là một danh từ dùng để chỉ hiện tượng siêu lạm phát.}

\vocabentry{Globalism}
{/ˈɡloʊbəlɪzəm/ (n)}
{The operation or planning of economic and foreign policy on a global basis.}
{Globalism has led to increased international trade but also concerns about the loss of national sovereignty.}
{Đây là một danh từ dùng để chỉ chủ nghĩa toàn cầu hóa kinh tế.}

\vocabentry{Outsourcing}
{/ˈaʊtsɔːrsɪŋ/ (n)}
{Obtain goods or a service from an outside or foreign supplier.}
{Outsourcing manufacturing to countries with lower labor costs can help companies reduce prices.}
{Đây là một danh từ dùng để chỉ việc thuê ngoài nhân lực hoặc sản xuất.}

\vocabentry{Offshoring}
{/ˌɔːfˈʃɔːrɪŋ/ (n)}
{The practice of basing some of a company's processes or services overseas.}
{Offshoring has been criticized for causing job losses in developed countries.}
{Đây là một danh từ dùng để chỉ việc dời các hoạt động kinh doanh ra nước ngoài.}

\vocabentry{Sustainability}
{/səˌsteɪnəˈbɪləti/ (n)}
{The ability to be maintained at a certain rate or level.}
{Economic sustainability requires balancing growth with environmental protection and social equity.}
{Đây là một danh từ dùng để chỉ sự phát triển bền vững.}

\vocabentry{Infrastructure}
{/ˈɪnfrəstrʌktʃər/ (n)}
{The basic physical and organizational structures needed for the operation of a society.}
{Investing in transport and communication infrastructure is essential for economic development.}
{Đây là một danh từ dùng để chỉ cơ sở hạ tầng kinh tế.}

\vocabentry{Human capital}
{/ˌhjuːmən ˈkæpɪtl/ (n)}
{The skills, knowledge, and experience possessed by an individual or population.}
{Education and training are the primary ways to increase a country's human capital.}
{Đây là một cụm danh từ dùng để chỉ vốn nhân lực.}

\vocabentry{Brain drain}
{/ˈbreɪn dreɪn/ (n)}
{The emigration of highly trained or intelligent people from a particular country.}
{Many developing countries suffer from a brain drain as their best doctors and engineers move abroad.}
{Đây là một cụm danh từ dùng để chỉ hiện tượng chảy máu chất xám.}

\vocabentry{Remittance}
{/rɪˈmɪtns/ (n)}
{A sum of money sent, especially by mail, in payment for goods or services or as a gift.}
{Remittances from workers abroad are a vital source of income for many families in poor countries.}
{Đây là một danh từ dùng để chỉ tiền kiều hối gửi về nước.}

\vocabentry{Globalization}
{/ˌɡloʊbələˈzeɪʃn/ (n)}
{The process by which businesses develop international influence or start operating on an international scale.}
{Globalization has led to a significant increase in cultural exchange and international cooperation.}
{Đây là một danh từ dùng để chỉ quá trình toàn cầu hóa.}

\vocabentry{Free trade}
{/ˌfriː ˈtreɪd/ (n)}
{International trade left to its natural course without tariffs, quotas, or other restrictions.}
{Free trade agreements aim to lower costs for consumers and increase competition between firms.}
{Đây là một cụm danh từ dùng để chỉ thương mại tự do.}

\vocabentry{Economic growth}
{/ˌekəˈnɑːmɪk ɡroʊθ/ (n)}
{An increase in the amount of goods and services produced per head of the population over a period of time.}
{The government's primary goal is to achieve sustainable economic growth and reduce poverty.}
{Đây là một cụm danh từ dùng để chỉ tăng trưởng kinh tế.}

\vocabentry{Inequality}
{/ˌɪnɪˈkwɑːləti/ (n)}
{Difference in size, degree, circumstances, etc.; lack of equality.}
{Rising income inequality is a major concern for many policymakers around the world.}
{Đây là một danh từ dùng để chỉ sự bất bình đẳng giàu nghèo.}

\vocabentry{Poverty line}
{/ˈpɑːvərti laɪn/ (n)}
{The estimated minimum level of income needed to secure the necessities of life.}
{Millions of people in the country still live below the poverty line.}
{Đây là một cụm danh từ dùng để chỉ ngưỡng nghèo.}

\vocabentry{Social welfare}
{/ˌsoʊʃl ˈwelfer/ (n)}
{Statutory procedure or social effort designed to promote the basic well-being of people in need.}
{The social welfare system provides a safety net for those who are unemployed or disabled.}
{Đây là một cụm danh từ dùng để chỉ phúc lợi xã hội.}

\vocabentry{Progressive tax}
{/prəˌɡresɪv ˈtæks/ (n)}
{A tax in which the tax rate increases as the taxable amount increases.}
{A progressive tax system is often used to redistribute wealth and reduce inequality.}
{Đây là một cụm danh từ dùng để chỉ thuế lũy tiến.}

\vocabentry{Universal basic income}
{/juːnɪˌvɜːrsl ˈbeɪsɪk ˈɪnkʌm/ (n)}
{A periodic cash payment delivered to all citizens on an individual basis without means test or work requirement.}
{Several countries are conducting trials of universal basic income to see if it can reduce poverty.}
{Đây là một cụm danh từ dùng để chỉ thu nhập cơ bản vô điều kiện.}

\vocabentry{Austerity}
{/ɔːˈsterəti/ (n)}
{Difficult economic conditions created by government measures to reduce a budget deficit.}
{The government's austerity measures have led to widespread protests and strikes.}
{Đây là một danh từ dùng để chỉ chính sách thắt lưng buộc bụng.}

\vocabentry{Stimulus}
{/ˈstɪmjələs/ (n)}
{A thing that rouses activity or energy in someone or something.}
{The economic stimulus package includes tax cuts and increased spending on infrastructure.}
{Trong kinh tế, đây là gói kích cầu hoặc kích thích kinh tế.}

\vocabentry{Regulation}
{/ˌreɡjuˈleɪʃn/ (n)}
{A rule or directive made and maintained by an authority.}
{Stronger financial regulations are needed to prevent another global economic crisis.}
{Đây là một danh từ dùng để chỉ quy định hoặc sự điều tiết của chính phủ.}

\vocabentry{Deregulation}
{/ˌdiːˌreɡjuˈleɪʃn/ (n)}
{The removal of regulations or restrictions, especially in a particular industry.}
{The deregulation of the airline industry led to lower fares and more choices for passengers.}
{Đây là một danh từ dùng để chỉ sự bãi bỏ quy định điều tiết.}

\vocabentry{Standard of living}
{/ˌstændərd əv ˈlɪvɪŋ/ (n)}
{The degree of wealth and material comfort available to a person or community.}
{The standard of living in the country has improved significantly over the last twenty years.}
{Đây là một cụm danh từ dùng để chỉ mức sống.}

\vocabentry{Purchasing power}
{/ˈpɜːrtʃəsɪŋ ˈpaʊər/ (n)}
{The financial ability to buy goods and services.}
{Inflation can significantly reduce the purchasing power of people on fixed incomes.}
{Đây là một cụm danh từ dùng để chỉ sức mua của đồng tiền.}

\vocabentry{Cost of living}
{/ˌkɔːst əv ˈlɪvɪŋ/ (n)}
{The amount of money needed to sustain a certain level of living.}
{The high cost of living in the city center is forcing many young families to move to the suburbs.}
{Đây là một cụm danh từ dùng để chỉ chi phí sinh hoạt.}

\vocabentry{Employment}
{/ɪmˈplɔɪmənt/ (n)}
{The state of having paid work.}
{The government is working on a program to boost employment in rural areas.}
{Đây là một danh từ dùng để chỉ việc làm.}

\vocabentry{Unemployment}
{/ˌʌnɪmˈplɔɪmənt/ (n)}
{The state of being unemployed.}
{High levels of youth unemployment can lead to social and political instability.}
{Đây là một danh từ dùng để chỉ tình trạng thất nghiệp.}

\vocabentry{Labor market}
{/ˈleɪbər ˈmɑːrkɪt/ (n)}
{The supply of people in a particular country or area who are able and willing to work in relation to the number of jobs available.}
{The labor market is becoming increasingly competitive for recent graduates.}
{Đây là một cụm danh từ dùng để chỉ thị trường lao động.}

\vocabentry{Productivity}
{/ˌproʊdʌkˈtɪvəti/ (n)}
{The effectiveness of productive effort, especially in industry.}
{Investing in new technology can significantly boost the company's productivity.}
{Đây là một danh từ dùng để chỉ năng suất lao động.}

\vocabentry{Output}
{/ˈaʊtpʊt/ (n)}
{The amount of something produced by a person, machine, or industry.}
{The factory's daily output has increased by 20\% since the new machines were installed.}
{Đây là một danh từ dùng để chỉ sản lượng đầu ra.}

\vocabentry{Input}
{/ˈɪnpʊt/ (n)}
{What is put in, taken in, or operated on by any process or system.}
{The cost of energy is a major input for the manufacturing sector.}
{Trong kinh tế là các yếu tố đầu vào.}

\vocabentry{Efficiency}
{/ɪˈfɪʃnsi/ (n)}
{The state or quality of being efficient.}
{We need to improve the efficiency of our production process to reduce costs.}
{Đây là một danh từ dùng để chỉ hiệu quả kinh tế hoặc năng suất.}

\vocabentry{Specialization}
{/ˌspeʃələˈzeɪʃn/ (n)}
{The process of concentrating on and becoming expert in a particular subject or skill.}
{International trade is based on the principle of regional specialization.}
{Đây là một danh từ dùng để chỉ sự chuyên môn hóa.}

\vocabentry{Comparative advantage}
{/kəmˈpærətɪv ədˈvæntɪdʒ/ (n)}
{The ability of an individual or group to carry out a particular economic activity more efficiently than another activity.}
{According to the theory of comparative advantage, countries should specialize in what they do best.}
{Đây là một cụm danh từ dùng để chỉ lợi thế so sánh.}

\vocabentry{Competitive advantage}
{/kəmˈpetətɪv ədˈvæntɪdʒ/ (n)}
{A condition or circumstance that puts a company in a favorable or superior business position.}
{A highly skilled workforce is a significant competitive advantage for any tech company.}
{Đây là một cụm danh từ dùng để chỉ lợi thế cạnh tranh.}

\vocabentry{Innovation}
{/ˌɪnəˈveɪʃn/ (n)}
{The action or process of innovating; a new method, idea, product, etc.}
{Continuous innovation is essential for staying ahead in the global market.}
{Đây là một danh từ dùng để chỉ sự đổi mới, sáng tạo.}

\vocabentry{Entrepreneurship}
{/ˌɑːntrəprəˈnɜːrʃɪp/ (n)}
{The activity of setting up a business or businesses, taking on financial risks in the hope of profit.}
{The university offers several programs to promote entrepreneurship among its students.}
{Đây là một danh từ dùng để chỉ tinh thần khởi nghiệp.}

\vocabentry{Startup}
{/ˈstɑːrtʌp/ (n)}
{A newly established business.}
{Many successful startups were founded in small garages or university dorm rooms.}
{Đây là một danh từ dùng để chỉ công ty khởi nghiệp.}

\vocabentry{Trade balance}
{/ˈtreɪd ˌbæləns/ (n)}
{The difference in value between a country's imports and exports.}
{The government is taking steps to improve the trade balance by boosting exports.}
{Đây là một cụm danh từ dùng để chỉ cán cân thương mại.}

\vocabentry{Current account}
{/ˌkɜːrənt əˈkaʊnt/ (n)}
{An account at a bank from which money may be withdrawn without notice.}
{The current account deficit has widened due to increased imports of energy.}
{Trong vĩ mô, đây là tài khoản vãng lai.}

\vocabentry{Capital account}
{/ˈkæpɪtl əˈkaʊnt/ (n)}
{An account used to record a company's investment in itself.}
{The capital account tracks the flow of investment money into and out of a country.}
{Trong vĩ mô, đây là tài khoản vốn.}

\vocabentry{Economic union}
{/ˌekəˈnɑːmɪk ˈjuːniən/ (n)}
{A type of trade bloc which is composed of a common market with a customs union.}
{The European Union is the world's most integrated economic union.}
{Đây là một cụm danh từ dùng để chỉ liên minh kinh tế.}

\vocabentry{Common market}
{/ˌkɑːmən ˈmɑːrkɪt/ (n)}
{A group of countries imposing few or no duties on trade with one another and a common tariff on trade with other countries.}
{The establishment of a common market allows for the free movement of goods, labor, and capital.}
{Đây là một cụm danh từ dùng để chỉ thị trường chung.}

\vocabentry{Customs union}
{/ˌkʌstəmz ˈjuːniən/ (n)}
{A group of countries that have agreed to charge the same import duties as each other and usually to allow free trade between themselves.}
{Member countries of the customs union have abolished all internal trade barriers.}
{Đây là một cụm danh từ dùng để chỉ liên minh thuế quan.}

\vocabentry{Free trade area}
{/ˌfriː ˈtreɪd ˈeriə/ (n)}
{The region encompassing a trade bloc whose member countries have signed a free trade agreement.}
{NAFTA created a massive free trade area between Canada, the US, and Mexico.}
{Đây là một cụm danh từ dùng để chỉ khu vực mậu dịch tự do.}

\vocabentry{Protection}
{/prəˈtekʃn/ (n)}
{Shielding a country's domestic industries from foreign competition.}
{Some countries use high tariffs as a form of protection for their local industries.}
{Trong thương mại là sự bảo hộ.}

\vocabentry{Sanitary}
{/ˈsænɪteri/ (adj)}
{Relating to the conditions that affect hygiene and health.}
{International trade is often subject to strict sanitary and phytosanitary regulations.}
{Trong thương mại là các rào cản kiểm dịch động thực vật.}

\vocabentry{Logistics}
{/ləˈdʒɪstɪks/ (n)}
{The detailed coordination of a complex operation involving many people, facilities, or supplies.}
{Efficient logistics are essential for businesses involved in international trade.}
{Đây là một danh từ dùng để chỉ ngành hậu cần/vận tải.}

\vocabentry{Supply chain}
{/ˈsʌplaɪ tʃeɪn/ (n)}
{The sequence of processes involved in the production and distribution of a commodity.}
{A disruption in the global supply chain can cause major shortages of goods.}
{Đây là một cụm danh từ dùng để chỉ chuỗi cung ứng.}

\vocabentry{Distribution}
{/ˌdɪstrɪˈbjuːʃn/ (n)}
{The action of sharing something out among a number of recipients.}
{The distribution of wealth in the country is highly unequal.}
{Trong kinh tế là sự phân phối hàng hóa hoặc thu nhập.}

\vocabentry{Wholesale}
{/ˈhoʊlseɪl/ (n/adj)}
{The selling of goods in large quantities to be retailed by others.}
{Wholesale prices have risen sharply due to increased transport costs.}
{Đây là một danh từ/tính từ dùng để chỉ việc bán sỉ/bán buôn.}

\vocabentry{Retail}
{/ˈriːteɪl/ (n/adj)}
{The sale of goods to the public in relatively small quantities for use or consumption.}
{The retail sector has been hit hard by the rise of online shopping.}
{Đây là một danh từ/tính từ dùng để chỉ việc bán lẻ.}

\vocabentry{Consumer}
{/kənˈsuːmər/ (n)}
{A person who purchases goods and services for personal use.}
{Consumer confidence is a key indicator of future economic activity.}
{Đây là một danh từ dùng để chỉ người tiêu dùng.}

\vocabentry{Producer}
{/prəˈduːsər/ (n)}
{A person, company, or country that makes, grows, or supplies goods or commodities for sale.}
{The country is one of the world's leading producers of coffee and cocoa.}
{Đây là một danh từ dùng để chỉ nhà sản xuất.}

\vocabentry{Merchant}
{/ˈmɜːrtʃənt/ (n)}
{A person or company involved in wholesale trade.}
{Venetian merchants grew rich through trade with the East during the Middle Ages.}
{Đây là một danh từ dùng để chỉ thương gia hoặc người buôn bán.}

\vocabentry{Trader}
{/ˈtreɪdər/ (n)}
{A person who buys and sells goods, currency, or shares.}
{He works as a currency trader for a large international bank.}
{Đây là một danh từ dùng để chỉ nhà giao dịch hoặc tiểu thương.}

\vocabentry{Broker}
{/ˈbroʊkər/ (n)}
{A person who buys and sells goods or assets for others.}
{You need a stock broker to help you buy and sell shares on the exchange.}
{Đây là một danh từ dùng để chỉ người môi giới.}

\vocabentry{Agent}
{/ˈeɪdʒənt/ (n)}
{A person who acts on behalf of another person or group.}
{She works as a real estate agent, helping people buy and sell houses.}
{Đây là một danh từ dùng để chỉ đại lý hoặc người đại diện.}

\vocabentry{Franchise}
{/ˈfræntʃaɪz/ (n)}
{An authorization granted by a government or company to an individual or group enabling them to carry out specified commercial activities.}
{He decided to open a fast-food franchise because of the proven business model.}
{Đây là một danh từ dùng để chỉ hình thức nhượng quyền kinh doanh.}

\vocabentry{Brand}
{/brænd/ (n)}
{A type of product manufactured by a particular company under a particular name.}
{The company has built a strong brand that is recognized all over the world.}
{Đây là một danh từ dùng để chỉ thương hiệu.}

\vocabentry{Marketing}
{/ˈmɑːrkɪtɪŋ/ (n)}
{The action or business of promoting and selling products or services.}
{Effective marketing is essential for the success of any new product.}
{Đây là một danh từ dùng để chỉ hoạt động tiếp thị.}

\vocabentry{Advertising}
{/ˈædvərtaɪzɪŋ/ (n)}
{The activity or profession of producing advertisements for commercial products or services.}
{The company spends millions of dollars on television and online advertising every year.}
{Đây là một danh từ dùng để chỉ hoạt động quảng cáo.}

\vocabentry{Promotion}
{/prəˈmoʊʃn/ (n)}
{The publicization of a product, organization, or venture so as to increase sales or public awareness.}
{The supermarket is running a special promotion on fresh fruit and vegetables.}
{Đây là một danh từ dùng để chỉ hoạt động khuyến mãi hoặc quảng bá.}

\vocabentry{Customer loyalty}
{/ˈkʌstəmər ˌlɔɪəlti/ (n)}
{The tendency of some consumers to continue buying the same brand of goods rather than competing brands.}
{Reward programs are a common way to build and maintain customer loyalty.}
{Đây là một cụm danh từ dùng để chỉ sự trung thành của khách hàng.}

\vocabentry{Market share}
{/ˈmɑːrkɪt ʃer/ (n)}
{The portion of a market controlled by a particular company or product.}
{The company is aggressively expanding its market share in Southeast Asia.}
{Đây là một cụm danh từ dùng để chỉ thị phần.}

\vocabentry{Market segment}
{/ˈmɑːrkɪt ˈseɡmənt/ (n)}
{A group of people who share one or more common characteristics.}
{The advertisement is aimed at a specific market segment: young professionals with high incomes.}
{Đây là một cụm danh từ dùng để chỉ phân khúc thị trường.}

\vocabentry{Market niche}
{/ˈmɑːrkɪt nɪtʃ/ (n)}
{A small, specialized section of a market for a particular product or service.}
{The company found a profitable market niche by selling organic baby clothes.}
{Đây là một cụm danh từ dùng để chỉ thị trường ngách.}

\vocabentry{Benchmarking}
{/ˈbentʃmɑːrkɪŋ/ (n)}
{Evaluate or check (something) by comparison with a standard.}
{Benchmarking our performance against our competitors helps us identify areas for improvement.}
{Đây là một danh từ dùng để chỉ việc đối chuẩn hoặc so sánh với tiêu chuẩn.}

\vocabentry{Best practice}
{/best ˈpræktɪs/ (n)}
{Commercial or professional procedures that are accepted or prescribed as being correct or most effective.}
{The company is adopting best practices in environmental management to reduce its carbon footprint.}
{Đây là một cụm danh từ dùng để chỉ quy trình tối ưu hoặc thông lệ tốt nhất.}

\vocabentry{Value added}
{/ˌvæljuː ˈædɪd/ (n)}
{The amount by which the value of an article is increased at each stage of its production.}
{High-tech manufacturing generates more value added than traditional assembly.}
{Đây là một cụm danh từ dùng để chỉ giá trị gia tăng.}

\vocabentry{Competitive}
{/kəmˈpetətɪv/ (adj)}
{Relating to or characterized by competition.}
{The company operates in a very competitive market with many large rivals.}
{Đây là một tính từ dùng để chỉ tính cạnh tranh.}

\vocabentry{Aggressive}
{/əˈɡresɪv/ (adj)}
{Ready or likely to attack or confront.}
{The company launched an aggressive marketing campaign to win back customers.}
{Trong kinh tế là chiến lược tấn công hoặc quyết liệt.}

\vocabentry{Dominant}
{/ˈdɑːmɪnənt/ (adj)}
{Most important, powerful, or influential.}
{The company has a dominant position in the global search engine market.}
{Trong kinh tế là vị thế thống trị thị trường.}

\vocabentry{Vibrant}
{/ˈvaɪbrənt/ (adj)}
{Full of energy and enthusiasm.}
{The city has a vibrant small business community and a booming tech sector.}
{Trong kinh tế là thị trường sôi động.}

\vocabentry{Stagnant}
{/ˈstæɡnənt/ (adj)}
{Showing no activity; dull and sluggish.}
{Economic growth has remained stagnant for several years due to a lack of investment.}
{Trong kinh tế là tình trạng trì trệ.}

\vocabentry{Voluntary}
{/ˈvɑːlənteri/ (adj)}
{Done, given, or acting of one's own free will.}
{The company made a voluntary commitment to reduce its greenhouse gas emissions.}
{Đây là một tính từ dùng để chỉ các thỏa thuận tự nguyện.}

\vocabentry{Mandatory}
{/ˈmændətɔːri/ (adj)}
{Required by law or rules; compulsory.}
{The new environmental standards are mandatory for all manufacturing companies.}
{Đây là một tính từ dùng để chỉ các quy định bắt buộc.}

\vocabentry{Legitimate}
{/ləˈdʒɪtɪmət/ (adj)}
{Conforming to the law or to rules.}
{The government aims to encourage legitimate business activity and discourage the shadow economy.}
{Đây là một tính từ dùng để chỉ tính hợp pháp hoặc chính đáng.}

\vocabentry{Illicit}
{/ɪˈlɪsɪt/ (adj)}
{Forbidden by law, rules, or custom.}
{Illicit trade in counterfeit goods is a major global problem.}
{Đây là một tính từ dùng để chỉ các hoạt động phi pháp.}

\vocabentry{Transparency}
{/trænsˈpærənsi/ (n)}
{The condition of being transparent.}
{Transparency in government spending is essential for preventing corruption.}
{Trong kinh tế là sự minh bạch tài chính.}

\vocabentry{Accountability}
{/əˌkaʊntəˈbɪləti/ (n)}
{The fact or condition of being accountable; responsibility.}
{There are calls for greater accountability for large financial institutions.}
{Đây là một danh từ dùng để chỉ trách nhiệm giải trình.}

\vocabentry{Integrity}
{/ɪnˈteɡrəti/ (n)}
{The quality of being honest and having strong moral principles.}
{The company is known for its high ethical standards and business integrity.}
{Trong kinh tế là sự chính trực trong kinh doanh.}

\vocabentry{Ethics}
{/ˈeθɪks/ (n)}
{Moral principles that govern a person's behavior or the conducting of an activity.}
{Business ethics are becoming increasingly important for environmentally conscious consumers.}
{Đây là một danh từ dùng để chỉ đạo đức kinh doanh.}

\vocabentry{Compliance}
{/kəmˈplaɪəns/ (n)}
{The action or fact of complying with a wish or command.}
{The bank has a dedicated team to ensure compliance with all financial regulations.}
{Đây là một danh từ dùng để chỉ sự tuân thủ các quy định pháp luật.}

\vocabentry{Bureaucracy}
{/bjʊˈrɑːkrəsi/ (n)}
{A system of government in which most of the important decisions are made by state officials.}
{Excessive government bureaucracy can discourage investment and hinder economic growth.}
{Trong kinh tế là bộ máy hành chính rườm rà.}

\vocabentry{Red tape}
{/ˌred ˈteɪp/ (n)}
{Excessive bureaucracy or adherence to rules and formalities, especially in public business.}
{Cutting red tape is essential for making the country more business-friendly.}
{Đây là một cụm danh từ chỉ các thủ tục hành chính quan liêu.}

\vocabentry{Bailout}
{/ˈbeɪlaʊt/ (n)}
{An act of giving financial assistance to a failing business or economy to save it from collapse.}
{Several large banks received a government bailout during the financial crisis.}
{Đây là một danh từ dùng để chỉ gói cứu trợ tài chính.}

\vocabentry{Safety net}
{/ˈseɪfti net/ (n)}
{A system of social welfare that protects people from the worst effects of poverty or unemployment.}
{A strong social safety net is essential for maintaining social stability during economic downturns.}
{Đây là một cụm danh từ dùng để chỉ mạng lưới an sinh xã hội.}

\vocabentry{Nationalization}
{/ˌnæʃnələˈzeɪʃn/ (n)}
{The transfer of a major branch of industry or commerce from private to state ownership or control.}
{The government announced the nationalization of the country's main oil company.}
{Đây là một danh từ dùng để chỉ sự quốc hữu hóa.}

\vocabentry{Subsidize}
{/ˈsʌbsɪdaɪz/ (v)}
{Support (an organization or activity) financially.}
{The state subsidizes public transport to keep it affordable for everyone.}
{Động từ của subsidy, chỉ việc trợ cấp.}

\vocabentry{Fiscal}
{/ˈfɪskl/ (adj)}
{Relating to government revenue, especially taxes.}
{The country's fiscal health has improved significantly thanks to higher tax revenues.}
{Đây là một tính từ dùng để chỉ thuộc về tài chính công hoặc ngân sách.}

\vocabentry{Economic}
{/ˌekəˈnɑːmɪk/ (adj)}
{Relating to economics or the economy.}
{The country is facing major economic challenges, including high inflation and unemployment.}
{Đây là một tính từ dùng để chỉ thuộc về kinh tế.}

\vocabentry{Economical}
{/ˌekəˈnɑːmɪkl/ (adj)}
{Giving good value or service in relation to the amount of money, time, or effort spent.}
{A small car is much more economical to run than a large SUV.}
{Đây là một tính từ dùng để chỉ sự tiết kiệm.}

\vocabentry{Wealth}
{/welθ/ (n)}
{An abundance of valuable possessions or money.}
{The discovery of natural gas has brought great wealth to the region.}
{Đây là một danh từ dùng để chỉ sự giàu có.}

\vocabentry{Prosperity}
{/prɑːˈsperəti/ (n)}
{The state of being prosperous.}
{The ultimate goal of all economic policy should be to increase national prosperity.}
{Đây là một danh từ dùng để chỉ sự thịnh vượng, thành công về kinh tế.}

\vocabentry{Standard}
{/ˈstændərd/ (n)}
{A level of quality or attainment.}
{We should work to improve the standard of living for all members of society.}
{Trong kinh tế thường dùng trong "standard of living" - mức sống.}

\vocabentry{Value}
{/ˈvæljuː/ (n/v)}
{The regard that something is held to deserve; the importance, worth, or usefulness of something.}
{The economic value of natural resources is often underestimated.}
{Đây là một danh từ/động từ dùng để chỉ giá trị hoặc việc định giá.}

